% !TeX program = pdflatex
\documentclass[11pt,a4paper]{article}
\usepackage[utf8]{inputenc}
\usepackage[T1]{fontenc}
\usepackage[brazilian]{babel}
\usepackage{amsmath,amssymb,amsthm}
\usepackage[margin=2.4cm]{geometry}
\usepackage{booktabs}
\usepackage{array}
\usepackage{enumitem}
\usepackage{xcolor}
\usepackage[hidelinks,breaklinks]{hyperref}
\usepackage[expansion=false]{microtype}

\sloppy
\emergencystretch=2em

\newtheorem{definicao}{Definição}
\newtheorem{teorema}{Teorema}
\newtheorem{proposicao}{Proposição}
\newtheorem{lema}{Lema}
\newtheorem{observacao}{Observação}
\newtheorem{corolario}{Corolário}

\newcommand{\code}[1]{\texttt{\small #1}}
\newcommand{\abs}[1]{\lvert #1\rvert}

\title{JEV e Campos de Contagem:\\
Construção de uma Camada de Representação para Decisão Tipada}
\author{}
\date{}

\begin{document}
\maketitle

\begin{abstract}
\noindent
Investiga-se se os contratos documentados de decisão tipada do JEV
(Choice, Score e Noul) admitem uma representação matemática sobre campos de
contagem induzidos por uma realização discreta
\(\pi\colon I\to X\), com
\(G(x)=\lvert\pi^{-1}(x)\rvert\) e
\(p_G(x)=G(x)/\lvert I\rvert\).

A distinção é central:
\[
\begin{array}{rcl}
\text{JEV} &\colon&
\text{estado}\to\text{pergunta tipada}\to\text{decisão},\\[2pt]
\text{construção deste paper} &\colon&
\text{realização}\to G\to\text{projeção }q\to G_q\to\text{distribuição}\to\text{leitura}.
\end{array}
\]
Não se afirma que o JEV implemente campos \(G\) internamente, nem que Tiffany
seja um ``JEV matemático'', nem que se tenha reconstruído o mecanismo do
modelo. Mostra-se apenas que Noul, Choice e Score são especializações de uma
mesma operação de projeção e normalização sobre \(G\), como
\emph{realizações matemáticas possíveis dos contratos de saída} no nível da
representação. A função \(q\) é um critério de classificação abstrato
(\(\text{question}\leadsto q\)), não a pergunta de API. A implementação
interna do JEV permanece indeterminada.
\end{abstract}

\tableofcontents
\newpage

% ═══════════════════════════════════════════════════════════════════════════
\section{Introdução}

O JEV (TypeSafe AI) é um modelo de decisão tipada: recebe um \emph{state} e um
mapa de perguntas tipadas e devolve respostas com probabilidades calibradas
conforme o tipo~\cite{typesafe-docs,jev-docs,systemone-primitives}. Os três
tipos documentados são:
\begin{itemize}[leftmargin=1.5em]
\item \textbf{Choice} --- selecção entre opções rotuladas, com distribuição e
confiança;
\item \textbf{Score} --- posição numa escala ordenada (2 a 10 níveis), com
distribuição, legenda e confiança; o valor pode ser fraccionário;
\item \textbf{Noul} --- probabilidade de \emph{yes} numa proposição binária
(sem campo de confiança separado).
\end{itemize}

Por outro lado, a arquitectura de campos de contagem (documentada em
\texttt{campos.tex}, com base em \texttt{fisica.tex}) formaliza
\[
I\xrightarrow{\pi}X\xrightarrow{\text{contagem}}G,
\qquad
G(x)=\lvert\pi^{-1}(x)\rvert,
\]
com conservação \(\sum_x G(x)=\lvert I\rvert\) e normalização
\(p_G(x)=G(x)/\lvert I\rvert\).

A pergunta deste paper é estritamente matemática:
\begin{quote}
É possível construir uma representação por campos de contagem que
\emph{exprime} os contratos Choice, Score e Noul, sem afirmar que o JEV
realize essa construção internamente?
\end{quote}

A resposta afirmativa limita-se ao nível do \emph{contrato de representação}.
A implementação, o treino e a calibração do JEV ficam fora do escopo.

% ═══════════════════════════════════════════════════════════════════════════
\section{O contrato JEV}

O que se utiliza abaixo é apenas o contrato público documentado
\cite{typesafe-docs,jev-docs,systemone-primitives,cloudflare-jev}.
\texttt{docs.typesafe.ai} é a documentação oficial da TypeSafe;
\texttt{jevtypesafeai.com} e \texttt{systemonemodels.org} são plataformas
independentes (não afiliadas à TypeSafe AI) que reproduzem o contrato; e
\texttt{developers.cloudflare.com} é uma página de integração de terceiros.
O estatuto aqui usado --- \emph{state}, questões tipadas, respostas
calibradas e os três primitivos --- é coerente entre todas elas.

\subsection{Estado e pergunta}

O pedido fundamental tem a forma
\[
(\textit{state},\;\{\textit{id}\mapsto\textit{question}\})
\;\longrightarrow\;
\{\textit{id}\mapsto\textit{answer}\}.
\]
O \emph{state} é o conteúdo a avaliar (texto ou estrutura serializável).
As perguntas são avaliadas em paralelo, cada uma em isolamento, sobre o mesmo
estado. O código da aplicação consome as respostas tipadas.

\subsection{Choice}

Uma pergunta Choice fornece \texttt{instructions} e um mapa
\texttt{criteria} de opções (até 255). A resposta documentada inclui:
\begin{itemize}[leftmargin=1.5em]
\item a opção seleccionada (\texttt{choice});
\item uma probabilidade para cada opção (\texttt{probabilities});
\item um valor de confiança (\texttt{confidence}) associado à opção
seleccionada.
\end{itemize}
Não se especifica aqui a fórmula interna de \texttt{confidence}; trata-se de
um campo documentado no contrato de saída, não de uma definição matemática
deste paper.

\subsection{Score}

Uma pergunta Score fornece \texttt{instructions} e uma lista ordenada de
níveis (2 a 10), do mais baixo ao mais alto. A resposta documentada inclui:
\begin{itemize}[leftmargin=1.5em]
\item um valor \texttt{score} (possivelmente fraccionário, entre níveis);
\item probabilidades por nível;
\item uma legenda dos níveis;
\item confiança.
\end{itemize}
A documentação descreve o score como uma posição na escala de níveis --- não a
medida de uma quantidade --- e define o valor fraccionário como a média
ponderada dos níveis pelas probabilidades (cada nível vezes a sua
probabilidade, somado sobre os níveis).

\subsection{Noul}

Uma pergunta Noul fornece \texttt{instructions} (e, opcionalmente, critérios
para \texttt{true}/\texttt{false}). A resposta é um único número
\(\texttt{noul}\in[0,1]\): a probabilidade calibrada de \emph{yes}. Não há
campo de confiança separado.

\subsection{Limite documental}

Nada no material público afirma que o JEV mantenha um campo de multiplicidade
\(G\), que conte realizações, ou que implemente Fourier/Mellin. Qualquer
construção deste paper que use \(G\) é, portanto, uma \emph{representação
externa} dos contratos, não uma descrição de implementação.

% ═══════════════════════════════════════════════════════════════════════════
\section{Realização e campo de contagem}

Reutiliza-se a notação de \texttt{campos.tex} / \texttt{fisica.tex}.

\begin{definicao}[Realização e multiplicidade]
\label{def:realizacao}
Seja \(I\) um conjunto finito (eventos / observações / índices de realização)
e \(X\) um conjunto discreto finito (células / estados observáveis). Uma
realização é uma aplicação \(\pi\colon I\to X\). O campo de multiplicidade
induzido é
\[
G(x)=\lvert\pi^{-1}(x)\rvert\in\mathbb N_0.
\]
\end{definicao}

\begin{proposicao}[Conservação]
\label{prop:conservacao}
\[
\sum_{x\in X}G(x)=\lvert I\rvert.
\]
\end{proposicao}

\begin{proof}
As fibras \(\{\pi^{-1}(x)\}_{x\in X}\) particionam \(I\). A aditividade da
cardinalidade sobre partições finitas dá a igualdade. (Resultado já
estabelecido na construção de base; cf.\ lemas de conservação em
\texttt{fisica.tex} e Proposição de conservação em \texttt{campos.tex}.)
\end{proof}

\begin{definicao}[Normalização]
\label{def:pg}
Quando \(\lvert I\rvert>0\),
\[
p_G(x)=\frac{G(x)}{\lvert I\rvert}.
\]
\end{definicao}

\begin{proposicao}
\label{prop:pg-soma}
\(\sum_x p_G(x)=1\).
\end{proposicao}

\begin{proof}
Imediato da Proposição~\ref{prop:conservacao}.
\end{proof}

\begin{observacao}[Estatuto de \(p_G\)]
\label{obs:pg-nao-jev}
\(p_G\) é a distribuição de frequência induzida pela realização \(\pi\).
\textbf{Não} é a probabilidade interna do JEV, nem uma regra de Born, nem um
output documentado da API. É objecto da construção Tiffany.
\end{observacao}

% ═══════════════════════════════════════════════════════════════════════════
\section{Projeção de um campo}

A operação central é agregar \(G\) através de uma classificação \(q\).

\begin{observacao}[Pergunta tipada \(\leadsto\) critério \(q\)]
\label{obs:question-q}
Uma pergunta JEV documentada contém \texttt{instructions} e, consoante o tipo,
\texttt{criteria} (opções, níveis ou true/false). Isso \textbf{não} é a função
matemática \(q\colon X\to Y\). A relação adoptada neste paper é apenas
estrutural:
\[
\boxed{\text{question}\;\leadsto\;\text{critério de classificação }q}
\]
e não \(\text{question}=q\). A função \(q\) é uma abstracção que fixa, sobre o
suporte \(X\), a partição relevante para agregar \(G\); a pergunta JEV é um
objecto de API com texto e critérios interpretados pelo modelo.
\end{observacao}

\begin{definicao}[Projeção / classificação]
\label{def:projecao}
Seja \(Y\) um conjunto finito e \(q\colon X\to Y\). O campo agregado é
\[
G_q(y)=\sum_{x\in q^{-1}(y)}G(x)=\lvert\pi^{-1}(q^{-1}(y))\rvert.
\]
\end{definicao}

\begin{proposicao}[Conservação sob projeção]
\label{prop:gq-conserva}
\[
\sum_{y\in Y}G_q(y)=\lvert I\rvert.
\]
\end{proposicao}

\begin{proof}
As pré-imagens \(q^{-1}(y)\) particionam \(X\); a soma dos \(G\) sobre essas
partes recupera \(\sum_x G(x)=\lvert I\rvert\).
\end{proof}

\begin{definicao}[Distribuição induzida]
\label{def:pq}
\[
p_q(y)=\frac{G_q(y)}{\lvert I\rvert}\qquad(\lvert I\rvert>0).
\]
(Pela Prop.~\ref{prop:gq-conserva}, \(\sum_y p_q(y)=1\).)
\end{definicao}

Os três casos de interesse são especializações de \(Y\):
\begin{itemize}[leftmargin=1.5em]
\item \(Y=\{0,1\}\) --- binário (Noul);
\item \(Y=C=\{c_1,\dots,c_n\}\) --- categorial (Choice);
\item \(Y=S=\{s_0,\dots,s_m\}\) ordenado --- níveis (Score).
\end{itemize}

% ═══════════════════════════════════════════════════════════════════════════
\section{Construção dos três contratos}

\subsection{Noul como caso binário}

\begin{definicao}[Proposição binária sobre \(X\)]
\label{def:P}
Seja \(P\colon X\to\{0,1\}\). Defina
\[
G_P(1)=\sum_{x\in X}G(x)\,P(x),\qquad
G_P(0)=\sum_{x\in X}G(x)\,(1-P(x)).
\]
\end{definicao}

\begin{proposicao}
\label{prop:noul-gq}
\(G_P(0)+G_P(1)=\lvert I\rvert\). Com \(\lvert I\rvert>0\),
\[
p_P(\mathrm{yes})=\frac{G_P(1)}{\lvert I\rvert},\qquad
p_P(\mathrm{no})=\frac{G_P(0)}{\lvert I\rvert},
\]
e \(p_P(\mathrm{yes})+p_P(\mathrm{no})=1\).
\end{proposicao}

\begin{proof}
Cálculo directo a partir das definições e da conservação de \(G\).
\end{proof}

\begin{observacao}[Comparação com o contrato Noul]
\label{obs:noul}
O contrato Noul devolve um único número em \([0,1]\): a probabilidade de
\emph{yes}. A construção acima produz \(p_P(\mathrm{yes})\) como frequência
normalizada de uma proposição \(P\) sobre o suporte da realização.
Classifica-se como
\[
\boxed{\text{realização possível do contrato Noul por campo de contagem.}}
\]
Não se afirma que o JEV calcule \(G_P\) nem que \(P\) seja a sua representação
interna do \emph{state}.
\end{observacao}

\subsection{Choice como caso categorial}

\begin{definicao}[Classificação categorial]
\label{def:q-choice}
Seja \(C=\{c_1,\dots,c_n\}\) um conjunto finito de opções e
\(q\colon X\to C\). Então
\[
G_q(c_i)=\sum_{x:\,q(x)=c_i}G(x),\qquad
p_q(c_i)=\frac{G_q(c_i)}{\lvert I\rvert}.
\]
\end{definicao}

\begin{proposicao}
\label{prop:choice-dist}
\(\sum_i p_q(c_i)=1\). Em particular, \(p_q\) é uma distribuição sobre as
opções.
\end{proposicao}

\begin{proof}
Proposição~\ref{prop:gq-conserva} e Definição~\ref{def:pq}.
\end{proof}

Uma leitura canónica da opção seleccionada, \emph{neste paper}, é
\[
c^\star\in\arg\max_{c\in C}p_q(c)
\]
(qualquer desempate fixo). Isto é uma política de leitura, não a política do
JEV.

\begin{definicao}[Concentração auxiliar]
\label{def:concentracao}
Para uso interno deste paper, define-se a medida de concentração
\[
\kappa(p_q)=\max_{c\in C}p_q(c)\in[0,1].
\]
\end{definicao}

\begin{observacao}
\label{obs:confidence}
\(\kappa\) \textbf{não} é o campo \texttt{confidence} do JEV. A fórmula
interna de confiança do JEV não é pública neste nível de detalhe e não é
reproduzida aqui. \(\kappa\) é apenas uma medida auxiliar de concentração da
distribuição \(p_q\).
\end{observacao}

\begin{observacao}[Comparação com o contrato Choice]
\label{obs:choice}
O contrato Choice devolve opção seleccionada, distribuição sobre opções e
confiança. A construção \(p_q\) fornece a distribuição; \(c^\star\) e \(\kappa\)
fornecem leituras possíveis de selecção e concentração.
Classifica-se como realização possível do contrato Choice por campo de
contagem, no nível da representação.
\end{observacao}

\subsection{Score como caso ordenado}

\begin{definicao}[Classificação ordenada]
\label{def:q-score}
Seja \(S=\{s_0,\dots,s_m\}\subset\mathbb R\), finito, com
\(s_0<\dots<s_m\) (tipicamente \(s_i=i\)), e \(q\colon X\to S\). Então
\[
G_q(s_i)=\sum_{x:\,q(x)=s_i}G(x),\qquad
p_q(s_i)=\frac{G_q(s_i)}{\lvert I\rvert}.
\]
\end{definicao}

\begin{definicao}[Valor esperado discreto]
\label{def:esperado}
\[
\mathbb E_G[S]=\sum_{i=0}^{m}s_i\,p_q(s_i).
\]
Como \(S\subset\mathbb R\), \(\mathbb E_G[S]\) é a média ponderada dos níveis
pela distribuição \(p_q\) (e pode cair entre dois níveis).
\end{definicao}

\begin{observacao}[Quatro objectos distintos]
\label{obs:score-distincao}
Não se confundem:
\begin{enumerate}[leftmargin=1.5em]
\item a distribuição \(\{p_q(s_i)\}\) sobre níveis;
\item o valor esperado \(\mathbb E_G[S]\) (possivelmente entre níveis);
\item um score discreto \(\arg\max_i p_q(s_i)\);
\item um score contínuo / fraccionário como funcional de leitura.
\end{enumerate}
O contrato Score do JEV documenta um valor (possivelmente fraccionário),
probabilidades por nível, legenda e confiança. A documentação oficial define o
score fraccionário exactamente como a média ponderada dos níveis --- o mesmo
funcional de \(\mathbb E_G[S]\). A coincidência é formal: nesta construção
\(\mathbb E_G[S]\) aplica o funcional à frequência \(p_q\) da realização; o
JEV aplica-o à distribuição calibrada produzida pelo modelo. Os dois inputs
não se identificam.
\end{observacao}

\begin{observacao}[Comparação com o contrato Score]
\label{obs:score}
A construção produz distribuição ordenada e um funcional de leitura
(\(\mathbb E_G[S]\) ou \(\arg\max\)). Classifica-se como realização possível
do contrato Score por campo de contagem. A interpretação de níveis como \emph{posições} numa escala, e não como medida
de uma quantidade física, está explícita na documentação (o score ``does not
measure the fraction of customers without a workaround'') e também não é
reivindicada aqui de modo diferente.
\end{observacao}

% ═══════════════════════════════════════════════════════════════════════════
\section{Unificação}

\begin{proposicao}[Forma unificada]
\label{prop:unificada}
Sejam \(Y\) finito e \(q\colon X\to Y\). A cadeia
\[
\boxed{
\begin{array}{c}
\text{realização }\pi
\;\longrightarrow\;
G
\;\longrightarrow\;
\text{projeção }q
\;\longrightarrow\;
G_q
\;\longrightarrow\;
p_q
\;\longrightarrow\;
\text{leitura }
L(p_q)
\end{array}
}
\]
especializa-se em:
\begin{itemize}[leftmargin=1.5em]
\item \(Y=\{0,1\}\) --- Noul (leitura: \(p_q(1)\));
\item \(Y=C\) categorial --- Choice (leitura: \(\arg\max p_q\), opcionalmente \(\kappa\));
\item \(Y=S\) ordenado --- Score (leitura: \(\mathbb E_G[S]\) ou \(\arg\max\)).
\end{itemize}
\end{proposicao}

\begin{proof}
Definições~\ref{def:projecao}--\ref{def:pq} e as secções anteriores.
\end{proof}

Esta é a construção central do paper. Os três tipos de decisão tipada
aparecem como \emph{casos de \(Y\) e de \(L\)}, não como três teorias
independentes.

% ═══════════════════════════════════════════════════════════════════════════
\section{Projeção conjunta e marginalização}
\label{sec:conjunta}

Uma realização admite várias classificações \(q_i\) sobre o mesmo campo
(Obs.~\ref{obs:multi-q}). A versão conjunta consiste em agregar \(G\) através
de todas as classificações em simultâneo; a leitura de cada pergunta isolada
torna-se, então, uma \emph{marginalização} do campo agregado. É uma extensão
matemática da construção unificada --- \emph{uma realização, múltiplas leituras
tipadas} --- e não uma alteração do contrato JEV.

\begin{definicao}[Projeção conjunta]
\label{def:projecao-conjunta}
Sejam \(Y_1,\dots,Y_n\) conjuntos finitos e \(q_i\colon X\to Y_i\). A projeção
conjunta é
\[
q=(q_1,\dots,q_n)\colon X\to Y_1\times\cdots\times Y_n,
\qquad
q(x)=\bigl(q_1(x),\dots,q_n(x)\bigr).
\]
O campo conjunto é
\[
G_q(y_1,\dots,y_n)
=
\sum_{\substack{x\in X\\ q_i(x)=y_i\;\forall i}}G(x)
=
\bigl\lvert\pi^{-1}\bigl(q^{-1}(y_1,\dots,y_n)\bigr)\bigr\rvert.
\]
Para \(n=1\), \(G_q\) coincide com a projeção simples
(Def.~\ref{def:projecao}): a definição conjunta é uma generalização.
\end{definicao}

\begin{proposicao}[Conservação conjunta]
\label{prop:conjunta-conserva}
\[
\sum_{(y_1,\dots,y_n)\in Y_1\times\cdots\times Y_n}G_q(y_1,\dots,y_n)=\lvert I\rvert.
\]
\end{proposicao}

\begin{proof}
As pré-imagens conjuntas \(q^{-1}(y_1,\dots,y_n)\) particionam \(X\) (\(q\) é
função); a soma recupera \(\sum_x G(x)=\lvert I\rvert\)
(Prop.~\ref{prop:conservacao}).
\end{proof}

\begin{definicao}[Normalização conjunta]
\label{def:pq-conjunta}
Quando \(\lvert I\rvert>0\),
\[
p_q(y_1,\dots,y_n)=\frac{G_q(y_1,\dots,y_n)}{\lvert I\rvert}.
\]
Pela Prop.~\ref{prop:conjunta-conserva}, \(\sum p_q=1\).
\end{definicao}

\begin{teorema}[Marginalização]
\label{thm:marginal}
Para cada \(i\in\{1,\dots,n\}\) e cada \(y_i\in Y_i\),
\[
\boxed{
\sum_{\substack{(y_1,\dots,y_n)\\ j\neq i}}G_q(y_1,\dots,y_n)
=
G_{q_i}(y_i),
}
\]
sendo a soma sobre todas as coordenadas excepto \(i\). Equivalentemente,
\[
\sum_{\substack{(y_1,\dots,y_n)\\ j\neq i}}p_q(y_1,\dots,y_n)
=
p_{q_i}(y_i).
\]
Cada projeção individual \(G_{q_i}\) (Def.~\ref{def:projecao}) é a
marginalização do campo conjunto sobre as restantes coordenadas.
\end{teorema}

\begin{proof}
Fixa \(y_i\in Y_i\). As fibras conjuntas com \(i\)-ésima coordenada \(y_i\)
particionam a pré-imagem \(q_i^{-1}(y_i)\):
\[
q_i^{-1}(y_i)
=
\bigsqcup_{(y_j)_{j\neq i}}
\bigl\{x\in X:\ q_j(x)=y_j\ \forall j\bigr\}.
\]
Somando \(G\) sobre essa partição e usando a definição de \(G_q\),
\[
\sum_{\substack{(y_1,\dots,y_n)\\ j\neq i}}G_q(y_1,\dots,y_n)
=
\sum_{x\in q_i^{-1}(y_i)}G(x)
=
G_{q_i}(y_i).
\]
A versão com \(p_q\) é a mesma soma dividida por \(\lvert I\rvert\)
(Def.~\ref{def:pq-conjunta}).
\end{proof}

\begin{corolario}[Os três contratos admitem projeções individuais e marginais conjuntas]
\label{cor:marginais}
Em primeiro lugar, cada contrato \textbf{já} se realiza por uma projeção
\emph{individual} \(G\to q_i\to G_{q_i}\) (Def.~\ref{def:projecao};
Obs.~\ref{obs:noul}, Obs.~\ref{obs:choice}, Obs.~\ref{obs:score}):
\begin{itemize}[leftmargin=1.5em]
\item coordenada binária \(Y_i=\{0,1\}\) --- Noul: \(p_P(\mathrm{yes})=p_{q_i}(1)\);
\item coordenada categorial \(Y_i=C\) --- Choice: distribuição \(p_{q_i}\),
\(c^\star=\arg\max\), \(\kappa\) auxiliar;
\item coordenada ordenada \(Y_i=S\) --- Score: distribuição \(p_{q_i}\) e
leitura \(\mathbb E_G[S]\) ou \(\arg\max\).
\end{itemize}
A marginalização não é, portanto, requisito para nenhum dos três. Quando várias
projeções \(q_1,\dots,q_n\) são consideradas \emph{simultaneamente} sobre o
mesmo \(G\), as suas distribuições individuais \(\{G_{q_i}\}\) são marginais do
campo conjunto \(G_q\) (Teor.~\ref{thm:marginal}):
\[
\boxed{
\text{individualmente, cada contrato basta-se com }G\to q_i\to G_{q_i};
}
\qquad
\boxed{
\text{em conjunto, }G_{q_i}=\sum_{j\neq i}G_q.
}
\]
\end{corolario}

\begin{proof}
Teor.~\ref{thm:marginal} aplicado a cada coordenada; as leituras são as das
Obs.~\ref{obs:noul}, Obs.~\ref{obs:choice} e Obs.~\ref{obs:score}.
\end{proof}

\begin{observacao}[Várias perguntas, uma realização]
\label{obs:multi-q-marg}
O contrato JEV avalia várias perguntas em paralelo sobre o mesmo
\textit{state} (Obs.~\ref{obs:multi-q}). Na construção, várias classificações
\(q_i\) sobre o mesmo \(G\) são coordenadas de uma única projeção conjunta;
cada saída tipada é a leitura de um marginal. Isto é uma propriedade da
construção --- \emph{uma realização, múltiplas leituras tipadas} --- e
\textbf{não} uma alegação de que o JEV implemente um campo conjunto \(G_q\).
\end{observacao}

\begin{observacao}[Objecto conceptual, leituras operacionais]
\label{obs:conjunta-itens}
O campo conjunto \(G_q\) vive no produto \(Y_1\times\cdots\times Y_n\), que
cresce com o número de perguntas. Mas as leituras expostas ao contrato são
\emph{marginais}: cada \(G_{q_i}\) calcula-se directamente por projeção simples
(Def.~\ref{def:projecao}), sem materializar o produto. O campo conjunto é a
descrição estrutural; os marginais são as leituras operacionais. Mantém-se a
separação entre representação construída e implementação (do JEV ou de um
pipeline Tiffany).
\end{observacao}

% ═══════════════════════════════════════════════════════════════════════════
\section{Memória no ambiente e leitura local}

Em \texttt{campos.tex}, o campo \(G\) é interpretado como traço ambiental
formal da realização: cada visita incrementa a ocupação da célula. A leitura
local (vizinhança, \(\arg\min G\), etc.) coordena agentes sem que nenhum
carregue a história completa.

\begin{observacao}
\label{obs:memoria}
A estrutura
\[
\text{acção}\to\text{traço em }G\to\text{leitura posterior}\to\text{nova acção}
\]
é compatível com o padrão de estigmergia discutido em \texttt{campos.tex}.
Não se afirma que isso prove estigmergia em geral, nem que o JEV opere por
estigmergia. Apenas se regista que \(G\) funciona, nesta construção, como
memória agregada da realização sobre a qual se definem projeções \(q\).
\end{observacao}

% ═══════════════════════════════════════════════════════════════════════════
\section{Campo e mudança de representação}

A arquitectura já estabelecida em \texttt{campos.tex} / \texttt{transformada.tex}
distingue:
\[
\boxed{\text{Fourier lê o incremento (coordenada aditiva);}}
\qquad
\boxed{\text{Mellin lê o salto (coordenada multiplicativa, via }u=\log x\text{).}}
\]
Magnitude e fase são componentes de uma decomposição polar de uma
representação complexa; Fourier e Mellin são transformadas associadas a
estruturas de grupo distintas. Não se identificam.

Quando o suporte admite estrutura de grupo (ou de escala com pesos
\(w_{k+1}=\rho_k w_k\)), pode aplicar-se uma transformação \(\mathcal T\) a
\(G\) ou a \(p_G\) \emph{depois} da normalização:
\[
G\;\longrightarrow\;p_G\;\longrightarrow\;\mathcal T(p_G).
\]
A projeção \(q\) que produz Choice/Score/Noul actua tipicamente sobre o
suporte \(X\) \emph{antes} ou \emph{independentemente} dessa mudança de
representação. Não se introduz relação entre \(\rho_k\) e \(p_q\).

No formalismo Walsh de \texttt{fisica.tex} (\(\chi_k\in\{\pm1\}\)), a
magnitude de \(\mathcal F(G)\) é determinada por \(G\); a base só transporta
sinal. Esse resultado não se exporta automaticamente para uma DFT complexa
sobre \(\psi=A e^{i\theta}\) sem hipóteses adicionais
(\texttt{campos.tex}, secção espectral).

% ═══════════════════════════════════════════════════════════════════════════
\section{Estado, campo e decisão}

Formaliza-se a diferença entre campo e decisão.

\begin{definicao}[Cadeia de representação e política]
\label{def:cadeia}
\[
\begin{array}{c}
\text{realização }\pi
\;\longrightarrow\;
G
\;\longrightarrow\;
\text{projeção }q
\;\longrightarrow\;
p_q
\;\longrightarrow\;
\text{leitura }L
\;\longrightarrow\;
\text{política }\Pi
\;\longrightarrow\;
\text{acção.}
\end{array}
\]
\end{definicao}

\begin{observacao}[Correspondência de contratos]
\label{obs:correspondencia}
Comparando com o fluxo documentado do JEV
\[
\textit{state}
\;\longrightarrow\;
\textit{question}
\;\longrightarrow\;
\textit{typed answer}
\;\longrightarrow\;
\textit{application code / action},
\]
obtém-se uma correspondência \emph{estrutural} de camadas:
\begin{center}
\small
\begin{tabular}{@{}lll@{}}
\toprule
Camada Tiffany & Papel & Camada JEV (contrato) \\
\midrule
realização / \(G\) & representação agregada & (não documentada) \\
\(q\) & critério de classificação
  (analogia; Obs.~\ref{obs:question-q}) & pergunta tipada (instruções + critérios) \\
\(p_q\) & distribuição de frequência & \texttt{probabilities} / \texttt{noul} \\
\(L(p_q)\) & leitura tipada & \texttt{choice} / \texttt{score} / \texttt{noul} \\
\(\Pi\) & política externa & código da aplicação \\
\bottomrule
\end{tabular}
\end{center}
A correspondência \textbf{não} implica identidade de implementação. O
\emph{state} do JEV não é identificado com \(G\); a pergunta tipada
\textbf{não} se identifica com a função \(q\) (Obs.~\ref{obs:question-q}).
O mesmo contrato de saída não implica o mesmo mecanismo interno.
\end{observacao}

\begin{observacao}[Várias projeções sobre o mesmo campo]
\label{obs:multi-q}
O contrato JEV permite várias perguntas em paralelo sobre o mesmo
\textit{state}. Na construção deste paper, o mesmo \(G\) admite várias
classificações \(q_1,\dots,q_n\), induzindo \(p_{q_1},\dots,p_{q_n}\) sem
alterar a realização. Isto é uma propriedade formal da cadeia
\(\pi\to G\to\{q_i\}\); não se afirma que o JEV implemente múltiplas
perguntas via um campo de contagem.
\end{observacao}

\subsection{O que a construção explica}

\begin{itemize}[leftmargin=1.5em]
\item Como obter distribuições normalizadas sobre \(\{0,1\}\), categorias ou
níveis ordenados a partir de uma realização e de uma classificação \(q\).
\item Como unificar os três contratos como casos de \(Y\) e de \(L\).
\item Como ler várias perguntas sobre a mesma realização: projeção conjunta
e marginalização (Sec.~\ref{sec:conjunta}).
\item Como separar representação (\(G\), \(p_q\)) de decisão (política
externa sobre a leitura).
\end{itemize}

\subsection{O que permanece externo}

\begin{itemize}[leftmargin=1.5em]
\item Como o JEV obtém \(p_q\) a partir do \emph{state} (modelo, treino,
calibração).
\item A fórmula de \texttt{confidence}.
\item Qualquer mecanismo interno de contagem ou de campo \(G\).
\end{itemize}

\subsection{Limites da correspondência}

A construção responde à pergunta de \emph{representabilidade} dos contratos.
Não responde à pergunta de \emph{como o JEV calcula}. Confundir as duas é
precisamente o que este paper evita.

% ═══════════════════════════════════════════════════════════════════════════
\section{Comparação estrutural JEV \(\times\) Tiffany}

\begin{center}
\small
\begin{tabular}{@{}p{0.22\textwidth}p{0.32\textwidth}p{0.38\textwidth}@{}}
\toprule
Camada & JEV (documentado) & Tiffany / este paper \\
\midrule
Entrada & \textit{state} & realização \(\pi\colon I\to X\) \\
Estrutura intermédia & (não especificada) & campo \(G\), \(p_G\) \\
Operação & julgamento tipado & contagem + projeção \(q\) \\
Espaço de saída & Choice / Score / Noul & \(p_q\) + leitura \(L\) \\
Probabilidade & \texttt{probabilities} / \texttt{noul} & \(p_q=G_q/\lvert I\rvert\) (frequência) \\
Transformação & estado \(\to\) decisão & \(\pi\to G\to\mathcal T(G)\) (opcional) \\
Escala & Score com níveis ordenados & pesos \(w_k\), razões \(\rho_k\) (outro eixo) \\
Coordenação & resposta consumível por código & leitura local / transformada \\
Memória & no \textit{state} fornecido & no campo \(G\) (traço) \\
\bottomrule
\end{tabular}
\end{center}

Formulação precisa:
\[
\boxed{\text{JEV formaliza a decisão;}\quad
\text{Tiffany formaliza a representação que antecede a leitura.}}
\]
Não:
``Tiffany é um JEV matemático.''

% ═══════════════════════════════════════════════════════════════════════════
\section{Tabela de estatuto}

\begin{center}
\small
\begin{tabular}{@{}p{0.18\textwidth}p{0.28\textwidth}p{0.28\textwidth}p{0.18\textwidth}@{}}
\toprule
Conceito & JEV (documentado) & Construção deste paper & Estatuto \\
\midrule
State & entrada da avaliação & --- (não identificado com \(G\)) & contrato \\
Question & Choice / Score / Noul (instruções + critérios) & \(q\) critério de classificação (\(\leadsto\), não \(=\)) & analogia estrutural \\
Choice & opção + probs + confidence & \(p_q\) sobre \(C\), \(\arg\max\), \(\kappa\) auxiliar & construção \\
Score & score (frac.) + probs + legend & \(p_q\) sobre \(S\), \(\mathbb E_G[S]\) (\(\neq\) fórmula interna) & construção \\
Noul & \(P(\mathrm{yes})\in[0,1]\) & \(p_P(\mathrm{yes})=G_P(1)/\lvert I\rvert\) & construção \\
Probability & saída calibrada do modelo & frequência \(p_G\), \(p_q\) & definição / cálculo \\
Confidence & campo documentado (Choice/Score) & \(\kappa=\max p_q\) (auxiliar; \textbf{não} confidence JEV) & definição auxiliar \\
Field \(G\) & \textbf{não documentado} & \(G(x)=\lvert\pi^{-1}(x)\rvert\) & definição \\
Projection \(q\) & \textbf{não documentado} & \(G_q(y)=\lvert\pi^{-1}(q^{-1}(y))\rvert\) & definição \\
Projeção conjunta \(q=(q_1,\dots,q_n)\) & \textbf{não documentado} & \(G_q(y_1,\dots,y_n)=\sum G(x)\) sobre \(q_i(x){=}y_i\) & definição (Def.~\ref{def:projecao-conjunta}) \\
Marginalização & --- & \(G_{q_i}(y_i)=\sum_{j\neq i}G_q\) (Teor.~\ref{thm:marginal}) & teorema \\
Várias perguntas em paralelo & avaliadas sobre o mesmo \textit{state} & uma realização, múltiplas leituras tipadas (marginais) & construção (Obs.~\ref{obs:multi-q-marg}) \\
Normalization & --- & \(p_q=G_q/\lvert I\rvert\) & definição \\
Decision & typed answer & leitura \(L(p_q)\) & comparação estrutural \\
Action & código da aplicação & política \(\Pi\) externa & observação \\
Memory & no state & no campo \(G\) & interpretação (Tiffany) \\
Fourier & --- & leitura aditiva (instância clássica) & reformulação \\
Mellin & --- & leitura multiplicativa via \(\log\) & reformulação \\
Implementação JEV & --- & \textbf{não determinada} & lacuna \\
\bottomrule
\end{tabular}
\end{center}

% ═══════════════════════════════════════════════════════════════════════════
\section{Conclusão}

Demonstrou-se a cadeia
\[
\boxed{
\begin{array}{c}
\text{realização }\pi
\;\to\;
G
\;\to\;
\text{projeção }q
\;\to\;
G_q
\;\to\;
\text{normalização }p_q
\;\to\;
\text{leitura }L
\;\to\;
\text{decisão / política.}
\end{array}
}
\]
Noul, Choice e Score aparecem como especializações:
\[
\begin{array}{rcl}
\text{Noul} &=& \text{caso binário }(Y=\{0,1\}),\\
\text{Choice} &=& \text{caso categorial }(Y=C),\\
\text{Score} &=& \text{caso ordenado }(Y=S).
\end{array}
\]

A conclusão \textbf{não} é: ``o JEV é implementado por campos de contagem'',
nem ``descobrimos como o JEV funciona''. A arquitectura pública do modelo não
expõe o mecanismo interno; este paper não o reconstrói.

A conclusão é:
\begin{quote}
Os contratos documentados de saída (Choice, Score, Noul) admitem uma
\emph{realização matemática} por projeção e normalização de um campo de
contagem \(G\), desde que se mantenha a distinção entre essa representação
construída e a implementação interna do sistema.
\end{quote}

Em particular: mesmo contrato de saída \(\neq\) mesmo mecanismo interno.

A pergunta em aberto para trabalho futuro é se Choice, Score e Noul podem ser
expressos como operações sobre um campo \(G\) \emph{de forma operacionalmente
útil} (por exemplo, em pipelines que já produzem realizações e ocupações),
sem pressupor a arquitectura interna do JEV --- e se um único \(G\) com
várias projeções \(q_i\), lidas como marginais do campo conjunto
(Sec.~\ref{sec:conjunta}, Teor.~\ref{thm:marginal}), é um modelo útil de
``uma realização, muitas leituras tipadas''.

\begin{thebibliography}{9}

\bibitem{typesafe-docs}
TypeSafe AI,
\emph{Documentation --- Choice, Score, Noul and System One}
(fonte primária: documentação oficial),
\url{https://docs.typesafe.ai/},
acedido em 2026.

\bibitem{jev-docs}
JevTypeSafeAI.com,
\emph{API decide --- state, questions, typed answers}
(plataforma independente, não afiliada à TypeSafe AI),
\url{https://jevtypesafeai.com/docs},
acedido em 2026.

\bibitem{systemone-primitives}
System One Models,
\emph{Choice, Score and Noul: the three primitives}
(hub independente que agrega a documentação oficial),
\url{https://systemonemodels.org/guides/choice-score-noul/},
acedido em 2026.

\bibitem{cloudflare-jev}
Cloudflare AI,
\emph{typesafe/jev --- structured evaluation model}
(página de integração de terceiros),
\url{https://developers.cloudflare.com/ai/models/typesafe/jev/},
acedido em 2026.

\bibitem{campos}
A.~M.~Lopes (projecto Tiffany),
\emph{Campos de Contagem e Mudança de Representação}
(\texttt{campos.tex}), repositório do projecto, 2026.

\bibitem{fisica}
A.~M.~Lopes (projecto Tiffany),
\emph{Física / formalismo de realização e multiplicidade}
(\texttt{fisica.tex}), repositório do projecto.

\bibitem{transformada}
A.~M.~Lopes (projecto Tiffany),
\emph{Transformada}
(\texttt{transformada.tex}), repositório do projecto.

\end{thebibliography}

\end{document}
